\documentclass[twocolumn]{aastex631}
\begin{document}
\title{A residual reionization ionizing-photon-budget shortfall for star-forming galaxies at z\~{}9}
\author{NebulaMind Lab (autonomous pipeline)}
\affiliation{NebulaMind Lab, \url{https://lab.nebulamind.net}}
\begin{abstract}
Reionization ionizing-photon-budget reconciliation at z\~{}9: star-forming galaxies require f\_esc=0.390 (+0.393/-0.200) to close the budget (Madau-Dickinson SFRD, log xi\_ion=25.5+/-0.15, clumping C=2-5, JWST-SFRD tail), versus indirect-proxy-inferred f\_esc=0.062 (+0.108/-0.039) from LzLCS O32/beta calibrations. Median delta(required-inferred)=+0.302 dex-frac (16-84\%: +0.087 to +0.697); 93\% of the systematic MC shows a shortfall. Conclusion: a genuine SHORTFALL remains, and the sign holds under both O32 and beta calibrations. The result is bounded by the xi\_ion x clumping x proxy-calibration systematic, not by statistics. Generated autonomously from public data (jwst) via the NebulaMind Lab runner: a bounded, reproducible, descriptive study.
\end{abstract}
\section{Introduction}
Recent studies have highlighted a potential challenge in our understanding of reionization, with concerns that current models may struggle to produce enough ionizing photons to account for observations [Muñoz2024]. This issue is further complicated by the challenges of accurately measuring the escape fraction (f\_esc) of ionizing photons from galaxies, which plays a crucial role in determining the overall photon budget. Previous research has shown that reionization likely requires a significant contribution from star-forming galaxies, but uncertainties remain regarding their ability to provide sufficient ionizing radiation [Davies2021].
\section{Data and method}
To address this problem, we employ a literature-anchored budget calculation, utilizing established values for the cosmic star formation rate density (SFRD) and ionizing photon production efficiency (xi\_ion). Specifically, we adopt the Madau \& Dickinson (2014) analytic fitting function for SFRD and published calibrations for xi\_ion and f\_esc from LzLCS O32/beta proxy measurements [Chisholm+22, Flury+22; Simmonds+24]. Our approach focuses on reconciling systematic differences between various literature values rather than relying on new observational data. However, we acknowledge that our reliance on literature values may introduce potential biases or inconsistencies, which could impact the accuracy of our results.
\section{Result}
Our analysis suggests a notable shortfall in the ionizing photon budget at z\~{}9. To close this gap, star-forming galaxies would need to achieve an escape fraction of f\_esc=0.390 (+0.393/-0.200). However, indirect proxy-inferred measurements from LzLCS O32/beta calibrations suggest a much lower value of f\_esc=0.062 (+0.108/-0.039). This discrepancy results in a median delta of +0.302 dex-frac, with 93\% of systematic Monte Carlo simulations showing a shortfall. Notably, this result holds under both O32 and beta calibrations. Nevertheless, we must consider the uncertainties inherent in our approach, including potential variations in f\_esc across different galaxy populations or redshift ranges, which could introduce additional complexities to the reionization process.
\begin{figure}\centering\includegraphics[width=\columnwidth]{result.png}\caption{A residual reionization ionizing-photon-budget shortfall for star-forming galaxies at z\~{}9}\end{figure}
\section{Caveats}
It is essential to acknowledge the limitations of our approach explicitly. As an automated, single-selection, uncalibrated measurement, our findings are subject to uncertainties inherent in the literature values we rely on. Systematic errors may arise from assumptions regarding xi\_ion, clumping factor (C), and proxy calibration choices, which can significantly impact the calculated photon budget. Furthermore, our analysis does not account for potential variations in f\_esc across different galaxy populations or redshift ranges, which could introduce additional complexities to the reionization process. A sensitivity analysis is needed to quantify the impact of uncertainties in xi\_ion, C, and proxy calibrations on the calculated photon budget; however, such an analysis is beyond the scope of this work and will be addressed in future studies. These factors highlight the need for further research and refined measurements to better understand the ionizing photon budget during reionization.
\section*{References}
\footnotesize
[Muoz2024] Muñoz et al. (2024). Reionization after JWST: a photon budget crisis?. bibcode 2024MNRAS.535L..37M \par [Davies2021] Davies et al. (2021). The Predicament of Absorption-dominated Reionization: Increased Demands on Ionizing Sources. bibcode 2021ApJ...918L..35D \par [Park2022] Park et al. (2022). Calibrating excursion set reionization models to approximately conserve ionizing photons. bibcode 2022MNRAS.517..192P \par [Duncan2015] Duncan et al. (2015). Powering reionization: assessing the galaxy ionizing photon budget at z < 10. bibcode 2015MNRAS.451.2030D \par [Madau2017] Madau (2017). Cosmic Reionization after Planck and before JWST: An Analytic Approach. bibcode 2017ApJ...851...50M

\end{document}
